\section{The closure--reality boundary}\label{sec:anchor:closure-reality-boundary}

The final row of the catalog is not another internal normal form.  It
marks the boundary between formal closure-of-record and bare physical
realisability.  A closure-of-record (\Cref{def:closure-formation}) is a
formed closure with an audit record of its closure data.  Bare physical
realisability says only that some physical or operational system
instantiates a model in an ordinary sense.  The question is whether the
two notions can be identified; without further structure they cannot.
The running example of \Cref{sec:intro} is a formal closure-of-record:
its carrier, split pair and repairs are audited data, and whether a
physical system instantiates it is the separate question treated here.

We model a closure-of-record abstractly.  Let \(\mathcal C\) be a
declared class of carriers, each considered over a declared scope
\(\mathfrak S\) that records the targets, routes, instruments,
interfaces, residuals, budgets and nonclaims under discussion.  (In this
chapter a carrier is an element of \(\mathcal C\), such as a subset of
a record set in the finite models, not the underlying set of the
working vocabulary of \Cref{sec:apparatus}.)  Formal closure over
\(\mathfrak S\) is given by a declared operator
\(\mathrm{Cl}_{\mathfrak S}:\mathcal C\to\mathcal C\), and a carrier is
formally closed when it is a fixed point.  The naive
\emph{closure-equals-reality} (CER) principle would identify bare
physical realisability of \(C\) with \(C=\mathrm{Cl}_{\mathfrak S}(C)\).
Its direct direction says that physical realisability implies
closure-of-record, and its converse that closure-of-record implies
physical realisability.  The two no-go theorems below show that neither
direction is valid without further structure.  Their implications use
no closure axiom; the finite models satisfy all three closure axioms,
so neither direction follows from those axioms.

Three operators occur in this chapter: an arbitrary operator
\(\mathrm{Cl}_{\mathfrak S}\) in the two no-go theorems; the five-flag
saturation \(\operatorname{Sat}_5\), a finite formal model of zero
structural defect, in clauses (a), (c) and (e) of \Cref{prop:F53}; and
the inference reflector \(\ClSB\) of \PaperAOR\
\citep{Tsiokos2026AOR}, which reflects ledgers into inference-closed
ledgers while keeping their evidence unchanged, in clause (b) only.
Full accounting in that apparatus requires supplied witnesses in
addition to inference closure.  Neither construction identifies formal
closure with bare physical realisability.

\begin{definition}[Reflective closure operator]\label{def:reflective-closure}
This definition fixes two different operators: the retained five-flag saturation \(\operatorname{Sat}_5\), defined here, and the AOR inference reflector \(\ClSB\), recalled from \citet{Tsiokos2026AOR}.
In the five-flag encoding used in this paper (called the \emph{retained} encoding below: it retains a carrier's structural data and changes only its five flags), a carrier over a scope \(\mathfrak S\) is a pair
\(C=(s,(P_1(C),\ldots,P_5(C)))\) of structural data \(s\) in a fixed set
\(\mathcal D_{\mathfrak S}\) and five propositions, read as: declared records are
typed, cited routes are audited, residuals are statused, nonclaims are
registered, and meta-audit obligations are discharged; two carriers are equal
when their data are equal and their flags are pairwise equivalent.
Write \(\operatorname{Sat}_5(C)\) for the operator keeping the structural
data and replacing each flag by a true proposition, and
\(\Xi_5(C)=\{i\in\{1,\ldots,5\}:\neg P_i(C)\}\).
The retained operational predicate is defined by \(\Xi_5(C)=\varnothing\).
This is a predicate on the five flags, not a construction of witnesses.

The inference-closure operator \(\ClSB\) of \PaperAOR, called the AOR operator below, is a different operator. Only \(\operatorname{Sat}_5\), \(\Xi_5\) and the retained operational predicate are defined here; the AOR notions below are recalled from \citet{Tsiokos2026AOR} for clause (b) of \Cref{prop:F53}. It takes a ledger to its
least rule-closed extension, retaining the actual obligation meanings,
evidence carrier, validity relation and certificate decoder. The reflected ledger
contains the seed ledger (the unit of the reflection). Inference closure does not supply missing
certificates. Concretely, a ledger over an audit model is a set of facts of
four kinds about obligations \(o\): \(o\) is declared, \(o\) is registered, \(e\)
is a certificate for \(o\), and \(o\) is discharged. A ledger is \emph{sound}
when every recorded certificate is valid for its obligation and the
meaning of every discharged obligation holds,
\emph{inference-closed} when it is closed under the rules of \PaperAOR, and
\emph{fully discharged} when every declared obligation is discharged. For a kernel state \(S\), \(\operatorname{Accounted}(S)\)
means that every reachable record has a supplied status and payload satisfying
its original readout. The \emph{canonical reflected presentation} of \(S\) is \(\ClSB\) applied to the ledger that presents \(S\) over its audit model; clause (b) of \Cref{prop:F53} therefore concerns the AOR reflector, not \(\operatorname{Sat}_5\). The relation between full discharge of the canonical reflected presentation and \(\operatorname{Accounted}(S)\) is the imported equivalence stated in clause (b) of \Cref{prop:F53} \citep{Tsiokos2026AOR}. Kernel states, their reachable records, supplied
statuses and payloads, and the canonical reflected presentation are the
notions of the AOR kernel of \citet{Tsiokos2026AOR}; they are used only in
clause (b) of \Cref{prop:F53}, whose equivalence is imported. Only the second sentence of clause (b), whose ledger is the scalar-budget ledger in the proof of \Cref{prop:F53}, is checked here, with the inference rules of \citet{Tsiokos2026AOR} stated in its proof; the first sentence is read in the notation of \citet{Tsiokos2026AOR}. In this paper
\(\ClSB\) always denotes that inference reflector, and no identification
\(\operatorname{Sat}_5=\ClSB\) is assumed.
\end{definition}

The two no-go theorems below derive each failure from a hypothesis that
supplies it; their general content is the clause of
\Cref{thm:CER-converse-no-go} for an arbitrary extensive operator: for an
extensive operator (in particular a closure operator) on the subsets of a set,
there exists a bare-realisability predicate, stable under every reduct
satisfying \(\operatorname{Red}C\subseteq C\), for which both CER implications
fail if and only if the operator is not the identity.

The \emph{reduct-obstruction hypothesis} of the next theorem concerns a
declared reduct map \(\operatorname{Red}:\mathcal C\to\mathcal C\), which forgets
a closure-relevant datum \(d\), and consists of three items:
\begin{enumerate}[label=(\roman*),topsep=2pt,itemsep=1pt]
  \item a carrier \(C^{+}\) with \(\operatorname{BareReal}_{\mathfrak S}(C^{+})\);
  \item \(\operatorname{BareReal}_{\mathfrak S}(C)\Rightarrow
    \operatorname{BareReal}_{\mathfrak S}(\operatorname{Red}C)\) for every
    \(C\in\mathcal C\);
  \item \(\mathrm{Cl}_{\mathfrak S}(\operatorname{Red}C^{+})\ne\operatorname{Red}C^{+}\).
\end{enumerate}
In the intended reading, \(\operatorname{Red}C^{+}\) keeps a target claim whose
audit uses \(d\), and (iii) holds because closure must add \(d\) back. The
proof uses all three items.

\begin{theorem}[Direct CER no-go]\label{thm:CER-direct-no-go}
Let \(\mathfrak S\) be a declared scope, let \(\mathcal C\) be a
class of carriers, let
\(\operatorname{BareReal}_{\mathfrak S}:\mathcal C\to\mathrm{Prop}\)
be bare physical realisability, and let
\(\mathrm{Cl}_{\mathfrak S}:\mathcal C\to\mathcal C\) be any operator on
carriers, the formal closure over \(\mathfrak S\).
Assume the reduct-obstruction hypothesis. Then \(C:=\operatorname{Red}C^{+}\) satisfies
\[
  \operatorname{BareReal}_{\mathfrak S}(C)
  \quad\text{and}\quad
  C\ne \mathrm{Cl}_{\mathfrak S}(C),
\]
so the direct CER implication
\[
  \operatorname{BareReal}_{\mathfrak S}(C)
  \Longrightarrow
  C=\mathrm{Cl}_{\mathfrak S}(C)
\]
fails at \(C\); by definition it fails exactly when some bare-realised carrier is not fixed by \(\mathrm{Cl}_{\mathfrak S}\). Moreover, the hypothesis is satisfiable by a closure
operator: on the subsets of \(\{d,e,f\}\), with
\(\mathrm{Cl}_{\mathfrak S}(C)=C\cup\{d\}\) if \(e\in C\) and
\(\mathrm{Cl}_{\mathfrak S}(C)=C\) otherwise, \(\operatorname{Red}C=C\setminus\{d\}\)
and \(\operatorname{BareReal}_{\mathfrak S}(C):\Longleftrightarrow C\subseteq\{d,e\}\),
the operator \(\mathrm{Cl}_{\mathfrak S}\) is extensive, monotone and
idempotent, \(\operatorname{BareReal}_{\mathfrak S}\) is stable under
\(\operatorname{Red}\), items (i)--(iii) hold with \(C^{+}=\{d,e\}\), and the
direct implication fails at \(\{e\}\). For every operator
\(\mathrm{Cl}_{\mathfrak S}\ne\mathrm{id}\) on \(\mathcal P(X)\), extensive or not (in particular every closure operator other than the identity), choose \(C_0\subseteq X\)
with \(\mathrm{Cl}_{\mathfrak S}(C_0)\ne C_0\). Defining
\(\operatorname{BareReal}_0(C):\Longleftrightarrow C\subseteq C_0\) gives a predicate
stable under every reduct satisfying \(\operatorname{Red}C\subseteq C\), for which
the direct implication fails at \(C_0\) (proved at the end of the proof below).
\end{theorem}
\begin{proof}
Assume the scope \(\mathfrak S\), the predicate
\(\operatorname{BareReal}_{\mathfrak S}\), the operator
\(\mathrm{Cl}_{\mathfrak S}\), and the reduct-obstruction
hypothesis.  Let \(C:=\operatorname{Red}C^{+}\). By (i) and (ii),
\[
  \operatorname{BareReal}_{\mathfrak S}(C),
\]
and by (iii),
\[
  C\ne \mathrm{Cl}_{\mathfrak S}(C).
\]
In the intended reading, the retained claim has a native residual whose
required record was the forgotten datum \(d\), and closure adds it back.
This \(C\) is the displayed witness to failure of the direct CER
implication. The finite-model clause is checked in the paragraph that follows; for the clause on arbitrary nonidentity operators, \(\operatorname{Red}C\subseteq C\subseteq C_0\) gives \(\operatorname{BareReal}_0(\operatorname{Red}C)\), and \(C_0\) satisfies \(\operatorname{BareReal}_0\) while \(\mathrm{Cl}_{\mathfrak S}(C_0)\neq C_0\).
\end{proof}

\paragraph{A finite model.} The hypothesis is satisfiable. Let the carriers
be the subsets of a record set \(\{d,e,f\}\), where \(e\) is a target claim
whose audit uses the datum \(d\); let \(\mathrm{Cl}_{\mathfrak S}(C)\) add \(d\)
to \(C\) whenever \(e\in C\), and let \(\operatorname{Red}C=C\setminus\{d\}\). Let \(\operatorname{BareReal}_{\mathfrak S}(C)\) mean
\(C\subseteq\{d,e\}\), the records of a realised system; it is preserved by
reducts. With \(C^{+}=\{d,e\}\), items (i)--(iii) hold, and the witness of the
theorem is \(C=\{e\}\): it is bare-realised, and
\(\mathrm{Cl}_{\mathfrak S}(C)=\{d,e\}\ne C\). The theorem states what any
carrier satisfying the hypothesis implies; the model shows that such carriers
exist. The operator \(\mathrm{Cl}_{\mathfrak S}\) of the model is extensive,
monotone and idempotent, so the model also shows that the direct CER
implication does not follow from the axioms of a closure operator together
with stability of bare realisability under reducts.

\subsection*{Layer instantiations}\label{layer-inst:CER-direct}

\begingroup\small
\begin{description}
\item[\emph{Common Rule informed-consent record
(medical research ethics).}] \textit{Analogy.}
The carrier \(C\) is the clinical procedure as
performed; the scope \(\mathfrak S\) imposes the
Common Rule audit (protocol, IRB approval,
informed consent, adverse-event reporting).  The
closure-relevant datum \(d\) is the signed
consent form with its IRB version history.  The
procedure can be bare-realised (physically
performed) while \(d\) is missing, so the reduct
satisfies \(\operatorname{BareReal}_{\mathfrak S}\)
but is not a closure-of-record: the retained
claim ``subject enrolled under IRB-approved
protocol'' has a residual whose required record
is the missing consent \citep{DHHS2018CommonRule}.

\end{description}
\endgroup

The \emph{closure-instantiation-gap hypothesis} of the next theorem
consists of two items:
\begin{enumerate}[label=(\roman*),topsep=2pt,itemsep=1pt]
  \item a formal carrier \(C\) whose records, residual statuses, bridges
    and nonclaims are closed over \(\mathfrak S\), so that
    \(C=\mathrm{Cl}_{\mathfrak S}(C)\);
  \item a witness that \(C\) has no bare physical instantiation, so that
    \(\neg\operatorname{BareReal}_{\mathfrak S}(C)\).
\end{enumerate}
The proof takes both items from the hypothesis.

\begin{theorem}[Converse CER no-go]\label{thm:CER-converse-no-go}
Let \(\mathfrak S\), \(\mathcal C\),
\(\operatorname{BareReal}_{\mathfrak S}\), and
\(\mathrm{Cl}_{\mathfrak S}\) be as above.  The converse CER implication
\[
  C=\mathrm{Cl}_{\mathfrak S}(C)
  \Longrightarrow
  \operatorname{BareReal}_{\mathfrak S}(C)
\]
fails for some carrier if and only if some fixed point of
\(\mathrm{Cl}_{\mathfrak S}\) is not bare-realised. In the model of
\Cref{thm:CER-direct-no-go}, whose closure operator is extensive, monotone and
idempotent and whose bare realisability is stable under reducts, the
closure-instantiation-gap hypothesis holds at \(C=\{f\}\) and the converse
implication fails there, so neither CER implication follows from the axioms of
a closure operator together with stability of bare realisability under
reducts. More generally, let \(\mathrm{Cl}_{\mathfrak S}\) be any
operator on the subsets of a set \(X\), and choose
\(C_0\subseteq X\) with \(\mathrm{Cl}_{\mathfrak S}(C_0)\ne C_0\), which exists
whenever \(\mathrm{Cl}_{\mathfrak S}\ne\mathrm{id}\). For
the predicate \(B_0(C):\Leftrightarrow C\subseteq C_0\) (the predicate \(\operatorname{BareReal}_0\) of \Cref{thm:CER-direct-no-go}) in place of
\(\operatorname{BareReal}_{\mathfrak S}\), \(B_0\) is stable under every reduct with
\(\operatorname{Red}C\subseteq C\) and the direct implication fails at \(C_0\); no other property of \(\mathrm{Cl}_{\mathfrak S}\) is used. If \(\mathrm{Cl}_{\mathfrak S}\) is extensive (in particular a closure operator), then \(\mathrm{Cl}_{\mathfrak S}(X)=X\) and the converse implication fails at \(X\). For every extensive operator,
including the identity, the always-false predicate, which is stable under
every reduct, makes the converse implication fail at \(X\). For every operator on \(\mathcal P(X)\), extensive or not, some reduct-stable predicate makes the direct implication fail if
and only if \(\mathrm{Cl}_{\mathfrak S}\ne\mathrm{id}\), since for the identity every
carrier is fixed. Hence, for every extensive operator, both CER implications fail for some reduct-stable
bare-realisability predicate if and only if \(\mathrm{Cl}_{\mathfrak S}\ne\mathrm{id}\).
The closure-instantiation-gap hypothesis at a carrier \(C\) is, item by item,
the statement that \(C\) witnesses the failure of the converse implication;
the content beyond this reformulation is the finite model and the clause on
extensive operators. In particular, under the closure-instantiation-gap
hypothesis there exists a carrier \(C\in\mathcal C\) such that
\[
  C=\mathrm{Cl}_{\mathfrak S}(C)
  \quad\text{and}\quad
  \neg\operatorname{BareReal}_{\mathfrak S}(C).
\]
\end{theorem}
\begin{proof}
Assume the scope \(\mathfrak S\), the predicate
\(\operatorname{BareReal}_{\mathfrak S}\), the operator
\(\mathrm{Cl}_{\mathfrak S}\), and the closure-instantiation-gap
hypothesis.  The hypothesis supplies a formal carrier \(C\) whose
records, residual statuses, bridges, and nonclaims are internally
closed:
\[
  C=\mathrm{Cl}_{\mathfrak S}(C).
\]
It also supplies the missing-instantiation witness:
\[
  \neg\operatorname{BareReal}_{\mathfrak S}(C).
\]
The point is that formal closure only checks the declared audit data
inside \(\mathfrak S\).  It does not by itself construct an external
physical system realising \(C\).  Therefore the displayed carrier
\(C\) is closed over \(\mathfrak S\) but not bare-realised, which is
the desired witness to failure of the converse CER implication. For the
last sentence of the theorem: the converse implication fails at a carrier
\(C\) exactly when its hypothesis \(C=\mathrm{Cl}_{\mathfrak S}(C)\) holds and
its conclusion \(\operatorname{BareReal}_{\mathfrak S}(C)\) does not. The
clause on the finite model is checked in the paragraph that follows. For the
general clause, \(\operatorname{Red}C\subseteq C\subseteq C_0\) gives stability;
\(C_0\) is bare-realised and not closed, so the direct implication fails there;
extensivity gives \(\mathrm{Cl}_{\mathfrak S}(X)=X\), and \(X\not\subseteq C_0\),
since \(X=C_0\) would make \(C_0\) closed, so the converse implication fails
at \(X\).
For the identity case, \(\mathrm{id}\) fixes every carrier, so the direct implication holds for every predicate; \(\bot\) is stable under every reduct and fails at the fixed carrier \(X\).
\end{proof}

\paragraph{A finite model.} In the model after \Cref{thm:CER-direct-no-go},
the carrier \(C=\{f\}\) contains no claim whose audit needs another record, so
\(\mathrm{Cl}_{\mathfrak S}(C)=C\), and \(f\notin\{d,e\}\), so
\(\neg\operatorname{BareReal}_{\mathfrak S}(C)\). Both items of the hypothesis
hold, and \(C\) witnesses the failure of the converse implication; in the same
model the direct implication fails at \(\{e\}\). So neither CER direction follows
from the closure-operator axioms together with reduct stability.

\subsection*{Layer instantiations}\label{layer-inst:CER-converse}

\begingroup\small
\begin{description}
\item[\emph{Construction documents for a
never-constructed building (architectural
design).}] \textit{Analogy.}
The carrier \(C\) is the building considered as a
closure-of-record: sealed construction documents,
stamped structural calculations, life-safety
analysis, code-review approvals, and issued
permit.  Under \(\mathfrak S\) the
architectural-design audit is closed at the
construction-document set.  When the owner
declines to break
ground or funding lapses, no physical building
instantiates \(C\), so \(C=\mathrm{Cl}_{\mathfrak
S}(C)\) yet \(\neg\operatorname{BareReal}
_{\mathfrak S}(C)\)
\citep{AllenIano2019Fundamentals}.

\end{description}
\endgroup


\paragraph{Reading the statement.}
The positive fixed-point equivalence below is proved for the retained
five-flag saturation. The AOR inference reflector has a separate
accounting criterion. An inference-closed ledger can still have an unpaid
obligation. Therefore the fixed-point equivalence for the flags cannot be
imported as an equivalence between inference closure and payment.
The proposition keeps both statements at their actual scopes.

\begin{proposition}[Closure--Reality Boundary resolution]\label{prop:F53}
For the retained encoding and the AOR kernel of
\Cref{def:reflective-closure}, the following statements hold.
\begin{enumerate}[label=(\alph*),topsep=2pt,itemsep=2pt]
\item A carrier is fixed by \(\operatorname{Sat}_5\) if and only if
  \(\Xi_5(C)=\varnothing\), equivalently all five flags hold.
\item (The accounting equivalence is imported; the kernel notions
  are those of \citet{Tsiokos2026AOR}, and the counterexample
  below is checked here.) For an AOR kernel state \(S\), its canonical reflected
  presentation is fully discharged if and only if
  \(\operatorname{Accounted}(S)\). Inference closure alone does not imply
  full discharge: a sound, inference-closed ledger over the scalar-budget
  evidence model of \PaperAOR, with all records declared, leaves the obligation \(5\le2\) unpaid.
\item The retained operational predicate, defined by
  \(\Xi_5(C)=\varnothing\), is consequently equivalent to the fixed-point
  predicate for \(\operatorname{Sat}_5\), and by (d) it is the unique
  predicate \(B\) for which both CER implications hold with
  \(\mathrm{Cl}=\operatorname{Sat}_5\) and \(B\) in place of bare
  realisability. No implication to or from bare
  physical realisability is asserted, and no identification with
  AOR inference closure is made.
\item For any operator \(\mathrm{Cl}\) and predicate \(B\), both CER
  implications hold for every carrier exactly when
  \(B(C)\iff C=\mathrm{Cl}(C)\) for every \(C\). Thus the fixed-point
  predicate is the unique predicate satisfying both implications.
\item If the set \(\mathcal D_{\mathfrak S}\) of structural data is nonempty, let
  \(D(C)=\{i\in\{1,\dots,5\}:P_i(C)\}\). For every proper subset
  \(D_0\subsetneq\{1,\dots,5\}\), both CER implications fail for
  \(\operatorname{Sat}_5\) and the predicate \(D(C)\subseteq D_0\). That
  predicate is stable under every reduct \(\operatorname{Red}\) with
  \(D(\operatorname{Red}C)\subseteq D(C)\) for all \(C\).
\end{enumerate}
Clause (b)'s accounting equivalence is imported from the AOR
kernel; its counterexample is checked using that same substrate. The other
clauses are proved for the retained encoding or for an arbitrary operator.
\end{proposition}
\begin{proof}
For (a), saturation keeps the structural data and sets all five flags to true.
Equality with the original carrier holds exactly when every original flag is
true, by propositional extensionality. This is precisely
\(\Xi_5(C)=\varnothing\). Clause (c) is (a) read through the definition of the retained predicate, together with (d) for \(\mathrm{Cl}=\operatorname{Sat}_5\).

For (b), \PaperAOR\ proves that a canonical discharged fact is derivable
exactly when its original record has a supplied valid status and payload.
Applying this equivalence to every reachable record yields the stated full
accounting equivalence. For the counterexample, take the scalar-budget model
with meaning \((c,b)\mapsto c\le b\), valid evidence a natural slack
\(e\) satisfying \(c+e=b\), and no dependency edges. Its interpreted
ledger declares and registers every obligation, stores exactly valid
certificates, and discharges exactly true inequalities. The inference rules of \citet{Tsiokos2026AOR} are four: from \(o\) declared
infer \(o\) registered; from a certificate \(e\) for \(o\) infer \(o\)
discharged; from \(o\) discharged infer \(o\) registered; and from \(o\)
declared and a dependency edge \(o\to d\) infer \(d\) declared. The
interpreted ledger contains every declaration and registration; a stored
certificate \(e\) for \((c,b)\) is valid, so \(c+e=b\) and \(c\le b\), and the
discharge the second rule concludes is already present; the model has no
dependency edges. Hence every rule preserves the ledger, and it is inference
closed. The declared obligation \((5,2)\)
is not discharged since \(5\nleq2\). Thus even a sound inference fixed point
need not be fully discharged.

For (d), the two implications are exactly the two directions of the stated
equivalence. Predicates with the same truth values are equal. For (e), keeping
any carrier's structural data and clearing its flags gives \(D(C)=\varnothing\),
which lies below \(D_0\) but is not fixed. Its saturation is fixed and has
all five flags, so it does not lie below the proper subset \(D_0\).
A reduct removing discharged flags only shrinks \(D(C)\).

The generic no-go theorems do not depend on identifying the two operators.
The retained saturation explains a formal fixed-point predicate. The
AOR construction additionally tracks the original witnesses required for
accounting; a consumer policy may require accepted delivery beyond accounting.
No bare physical existence claim follows from either construction.
\end{proof}

\subsection*{Layer instantiations}\label{layer-inst:F53}

\begingroup\small
\begin{description}
\item[\emph{Good Manufacturing Practice
(pharmaceutical regulatory affairs).}] \textit{Analogy.}
The raw carrier is a manufactured drug substance and the declared scope
includes batch records, validated processes, traceability and quality-unit
release. Closing the documentation's inference rules does not establish
that a missing validation was performed. Actual certificates must satisfy
those original obligations; release may additionally require the consumer's
acceptance policy. The distinction illustrates inference closure versus
accounting, without deriving a regulatory guarantee \citep{ICH2016Q7}.

\end{description}
\endgroup


This chapter closes the catalog by separating formal fixed points, audited
accounting and bare physical realisability. Neither CER direction follows
from the closure-operator axioms with reduct stability. The retained positive
equivalence concerns \(\operatorname{Sat}_5\); the AOR reflector
requires the separate source-accounting condition. The chapter imports that
condition's kernel theorem and does not claim that inference closure creates
its witnesses. This chapter does not re-prove the AOR apparatus or extend its
source-accounting equivalence to non-AOR reflective frameworks. The broader philosophical setting is the structural-realism
debate launched by \citet{Worrall1989StructuralRealism}; this chapter records
the formal closure / realisability gap without taking a stand in that debate.
